
\chapter{JUSTIFICATIVA}
\label{capitulo-justificativa}
A realização desta atividade de extensão justifica-se pela oportunidade de aproximar conceitos de computação de ponta da realidade cotidiana da comunidade. Frequentemente, a infraestrutura em nuvem corporativa e a administração de sistemas baseados em Linux são vistas como temas complexos e de difícil absorção para pessoas fora do ambiente profissional de TI. Utilizar o ecossistema do \textit{Minecraft} como estudo de caso prático funciona como um estímulo de engajamento, permitindo que conceitos mais complexos de redes de computadores, roteamento e segurança sejam apreendidos a partir de um projeto simples e de utilidade imediata.

Sob a perspectiva econômica e técnica, o projeto valida uma alternativa de custo zero, cuja capacidade supera as soluções comerciais populares de entrada. Enquanto opções pagas no mercado brasileiro exigem mensalidades recorrentes por instâncias restritas a pouca memória e processadores limitados, a camada \textit{Always Free} da Oracle Cloud oferece, na arquitetura Ampere ARM, até 2 OCPUs e 12 GB de RAM. Com essas configurações disponibilizadas e a vinculação de um endereço IPv4 público, com latência baixa, devido à proximidade do local onde o servidor está localizado (São Paulo/SP), permitindo uma autonomia operacional integral.

Além da economia financeira direta para os usuários atendidos, a relevância da proposta também está na abrangência dos tópicos técnicos ensinados. O conhecimento adquirido não se restringe à operação do jogo. O participante é capacitado para administrar um nó Linux completo, habilitando-o a instalar aplicações corporativas ou de uso pessoal, como centrais de mídia (\textit{Plex ou Jellyfin}), servidores de correio eletrônico, servidores web de documentação ou até mesmo túneis de rede privada.

Se a ação não fosse concretizada, os usuários permaneceriam sujeitos a custos contínuos desnecessários ou restritos a ambientes de baixa confiabilidade, comprovando a visão de que a computação em nuvem é inacessível para fins não corporativos.